\lecture{15}
\title{Double Descent and Implicit Regularization}

\begin{document}

\subsection{Double Descent}

The intuition that classical statistics gives us
is that models with larger hypothesis classes $\mathcal{H}$
are more prone to overfitting.
In this sense, achieving zero training error
via increasing the parameter count
may have negative consequences on the test error.

\begin{center}
    \frame{\includegraphics[width=0.4\textwidth]{lecture-15/bias-variance}}
\end{center}

The modern deep learning perspective, however,
is that setting the parameter count \emph{even higher}
can contribute positively toward
decreasing the test error in the long term!
This phenomenon is known as \textbf{double descent}.

\begin{center}
    \frame{\includegraphics[width=0.4\textwidth]{lecture-15/double-descent}}
\end{center}

The graphs below provide some intuition
for why double descent takes place.

\begin{center}
    \frame{\includegraphics[width=0.75\textwidth]{lecture-15/double-descent-graphs}}
\end{center}

Here's what these graphs indicate:
\begin{itemize}
    \item In the leftmost graph,
    the number of parameters is low,
    so both the training and test error are large.
    \item In the middle graph,
    the number of parameters is just large enough
    for ERM to find a perfect interpolation
    through the training data---the
    \emph{interpolation threshold}.
    Thus, training error is zero,
    but test error is large due to overfitting.
    \item In the rightmost graph,
    the number of parameters is large enough
    for ERM to find \emph{multiple} interpolations
    through the training data.
    If we add a ridge penalty $\lambda \|\theta \|^2$
    to our loss function, then ERM finds the interpolation
    with the smallest weights, which overfits the least!
    So the training error is zero,
    and the test error is small.
\end{itemize}

\begin{remark}
    Besides plain intuition,
    one solid argument as to why adding a ridge penalty
    should improve test error
    is the error bound we derived in
    \href{/6.790/lectures/10/#estimation-error-bounding}{Lecture 10}
    (restated below).
    \begin{quote}
        Assuming $\mathcal{L}(x, y) \in [0, 1]$,
        for all $\delta \in (0, 1)$
        we have with probability $1 - \delta$ that:
        \[
        \mathcal{R}(\hat{h}) - \mathcal{R}(h^*)
        \leq 2 \cdot 
        \sqrt{\frac{\log |\mathcal{H}| + \log\left(2 / \delta\right)}{2n}}.
        \]
    \end{quote}
    The ridge penalty decreases the size $|\mathcal{H}|$
    of the set of reasonable hypotheses.
\end{remark}

Here's a particular demonstration of double descent
on the MNIST data set.

\begin{center}
    \frame{\includegraphics[width=0.4\textwidth]{lecture-15/MNIST}}
\end{center}

\subsection{Implicit Regularization: Regression}

This section aims to communicate the following \emph{vibe}:
\begin{quote}
    \textbf{Vibe.}
    The solution that gradient descent lands on
    will naturally tend to have small weights---even
    without a ridge penalty term!
\end{quote}
We justify this vibe
by demonstrating its validity specifically in the case of
linear and logistic regression.

\begin{theorem}{Linear Regression Implicit Regularization}
    Consider a linear regression
    $h \colon \mathbb{R}^d \to \mathbb{R}$
    given by $h(x) := w^{\top}x$
    with training error $\hat{\mathcal{R}}(w) := \|Xw - Y\|^2$.
    Furthermore, assume that $d > n$, meaning
    zero training error $\hat{\mathcal{R}}(w) = 0$ is achievable.

    Among all parameters $w$ for which $\hat{\mathcal{R}}(w) = 0$
    (or equivalently, $Xw = Y$),
    let $w^*$ denote the parameter with the lowest magnitude.
    Then GD with $w_0 = 0$ converges to $w^*$.
\end{theorem}

\begin{proof}
    As usual, let $X \in \mathbb{R}^{n \times d}$ and $Y \in \mathbb{R}^n$
    encode the training data.
    Suppose GD with $w_0 = 0$ converges to $w_{\infty}$.
    \begin{quote}
        \textbf{Lemma 1.}
        We have $w_{\infty} \in \mathrm{colspan}(X^{\top})$.

        \emph{Proof:}
        Both $w_0 = 0$ and every update $\nabla \hat{\mathcal{R}}(w_t)
        \propto X^{\top}(Xw_t - Y)$
        lie in $\mathrm{colspan}(X^{\top})$. ~ $\square$
    \end{quote}
    Now consider any parameter $w$ for which
    $\hat{\mathcal{R}}(w) = 0$ (equivalently, $Xw = Y$).
    Since $\hat{\mathcal{R}}(w_{\infty}) = 0$ as well,
    we also have $Xw_{\infty} = Y$.
    Therefore, if we denote $v := w - w_{\infty}$, we have:
    \[ Xv = Xw - Xw_{\infty} = Y - Y = 0
    ~ \implies ~
    v \in \mathrm{ker}(X) = \mathrm{colspan}(X^{\top})^{\perp}. \]
    In particular, since $w_{\infty} \in \mathrm{colspan}(X^{\top})$
    by Lemma 1, this implies $\langle w_{\infty}, v \rangle = 0$.
    Now the Pythagorean theorem gives
    $\|w\|^2 = \|w_{\infty}\|^2 + \|v\|^2 \geq \|w_{\infty}\|^2$,
    with equality precisely when $w = w_{\infty}$,
    as desired. ~ $\blacksquare$
\end{proof}

\begin{remark}
    The exact same proof works just as well for SGD;
    all that changes is that the proof of Lemma 1
    must specify that every possible update
    still lives in $\mathrm{colspan}(X^{\top})$,
    which is easy to verify.
\end{remark}

It turns out the same phenomenon
can be proven true for logistic regression as well;
proof omitted.

\begin{theorem}{Logistic Regression Implicit Regularization}
    Given a linearly separable dataset,
    let $\mathcal{W}$ denote the set of all parameters
    $w$ such that $y_i \langle w, x_i \rangle \geq 1$ for all $i$.
    Define $w^* := \argmin_{w \in \mathcal{W}} \|w\|^2$.

    Suppose we train a logistic regression on this dataset using GD
    with sufficiently small $\eta$ and any initialization $w_0$.  Then:
    \[ w_t = w^* \log t + O(\log \log t)
    ~ \text{ as } t \to \infty. \]
    In particular, this implies that
    $\lim_{t \to \infty} \frac{w_t}{\|w_t\|} = \frac{w^*}{\|w^*\|}$,
    so GD converges in the direction of $w^*$.
\end{theorem}

\subsection{Implicit Regularization: Early Stopping}

It has also been observed that stopping training early
can be a useful tactic for reducing testing error.
For example, one might proceed with SGD
until validation error increases.

\begin{center}
    \frame{\includegraphics[width=0.4\textwidth]{lecture-15/early-stopping}}
\end{center}

Informally, we can say the following;
justification omitted.

\begin{theorem}{Yao et al., 2007}
    For least squares regression,
    stopping GD early at time step $t$
    is equivalent to ridge penalty $\frac{1}{\eta t} \|w\|^2$.
\end{theorem}

Early stopping usually occurs
\emph{before} the interpolation threshold,
meaning we stop before double descent becomes relevant.

\end{document}